\documentclass[10pt,a4paper]{amsart}
\usepackage[margin=19mm]{geometry}
\usepackage{amsmath,amssymb,mathtools}
\usepackage[scr=rsfs,cal=euler]{mathalfa}
\usepackage{xfrac}
\usepackage[unicode,hidelinks,bookmarks=false]{hyperref}
\newcommand{\ie}{i.e.\ }
\newcommand{\cf}{cf.\ }
\newcommand{\resp}{respectively\ }
\newcommand{\Subsection}{\subsection}
\providecommand{\nobreakdash}{\nobreak\hbox{-}}
\setlength{\emergencystretch}{2em}
\multlinegap0pt
\usepackage{bm}
\usepackage{bbm}
\usepackage{dsfont}
\usepackage{xspace}
\usepackage{cancel}
\usepackage{mathtools}
\usepackage{amsthm}
\makeatletter
\@ifundefined{theorem}{\newtheorem{theorem}{Theorem}}{}
\@ifundefined{lemma}{\newtheorem{lemma}[theorem]{Lemma}}{}
\makeatother
\DeclareMathOperator{\rank}{rank}
\DeclareMathOperator{\spec}{spec}
\newcommand{\set}[2]{\left\{#1 \mid #2\right\}}
\title[A new perspective on Willems' fundamental lemma: Universalit]{A new perspective on Willems' fundamental lemma: Universality of persistently exciting inputs}
\author{Amir Shakouri, Henk J. van Waarde, M. Kanat Camlibel}
\date{Source: arXiv:2503.12489v2 (CC BY 4.0)}
\begin{document}
\maketitle

\noindent\emph{Source and licence.} A new perspective on Willems' fundamental lemma: Universality of persistently exciting inputs. Version arXiv:2503.12489v2 (CC BY 4.0), distributed under the Creative Commons Attribution 4.0 International licence, as recorded in the arXiv licence metadata for this exact version and shown on the abstract page https://arxiv.org/abs/2503.12489v2. Full rights notice: licenses/arxiv-2503.12489-CC-BY-4.0.txt.\par\smallskip

\noindent\textit{Editorial scope.} Counted unit: the Lemma whose source label is \texttt{lem:henk} in the article (source0.tex lines 293--308, source label \texttt{lem:henk}), delivered together with the article's own complete original proof of it (source0.tex lines 310--404, one \texttt{proof} environment of 95 lines). No mathematics is written, repaired or re-derived: statement and proof are the article's own text line for line. Required context retained uncounted: none. Every definition, display and label of those lines is retained unchanged. Omitted from this package and not redistributed: 1--292, 309--309, 405--668. Nothing inside a retained excerpt is omitted or altered: every retained body is delivered exactly as the article's lines are written, so no omission receipt of any kind accompanies this package. Citations: the retained excerpts carry no citation command at all, so no bibliography entry of the article is transcribed and no reference is redistributed. Reference adaptation: none was needed and none was made; every reference inside the delivered unit resolves inside it, so no unresolved reference or citation is produced. Printed slips kept literal: no printed word, symbol, hypothesis or number of the counted statement or of its proof was altered. Layout and class: the article's own class is \texttt{ieeecolor} with packages including mathdots, enumitem, generic, comment, mathtools, amsmath, tabularx, diagbox, graphicx, mathrsfs, cite, amssymb, amsfonts, algorithmic; the wrapper is the lane's pinned \texttt{amsart} editorial wrapper at 10pt on A4, which restates the theorem-like environments the retained text needs and adds the article's own macro definitions that the retained text needs (\texttt{\textbackslash rank}, \texttt{\textbackslash set}, \texttt{\textbackslash spec}), with extra wrapper packages (bm, bbm, dsfont, xspace, cancel, mathtools). The delivered numbering is the unit's own. Assets: no retained span contains a figure environment, a graphics-inclusion command, a PDF-page inclusion, a table or a TikZ picture; no image file is copied into this package, the package declares no asset dependency and no image file is redistributed with it. Licence: version arXiv:2503.12489v2 of this article is distributed under the Creative Commons Attribution 4.0 International licence, as recorded in the arXiv licence metadata for this exact version and shown on the abstract page https://arxiv.org/abs/2503.12489v2. A full rights notice is delivered at licenses/arxiv-2503.12489-CC-BY-4.0.txt. Characters: every retained excerpt is pure ASCII, so the delivered engine, XeLaTeX with its default Latin Modern text font, reports no missing character.
\begin{lemma}
\label{lem:henk}
Let $T\in\mathbb{N}$, $L\in[1,T]$, and $u_{[0,T-1]}\in\mathbb{R}^{mT}$. If $u_{[0,T-1]}$ is not persistently exciting of order $n+L$, then there exists a controllable pair $(A,B)$ and a state $x_{[0,T-L]}$ satisfying $\begin{bmatrix}
u_{[0,T-L]} \\ x_{[0,T-L]}
\end{bmatrix}\in\mathfrak{B}_{T-L+1}(A,B)$ such that
\begin{equation}
\label{eq:lem:henk-1}
\begin{bmatrix}
v^\top & w^\top
\end{bmatrix}\begin{bmatrix}
\mathcal{H}_L(u_{[0,T-1]}) \\
\mathcal{H}_1(x_{[0,T-L]}) 
\end{bmatrix}=0
\end{equation}
for some $v\in\mathbb{R}^{mL}$ and $w\in\mathbb{R}^n\backslash\{0\}$.
\end{lemma}

\begin{proof}
We distinguish two cases: $T<n+L-1$ and \mbox{$T\geq n+L-1$}. In case $T<n+L-1$, we have
\begin{equation}
\rank\mathcal{H}_1(x_{[0,T-L]})\leq \min\{n,T-L+1\}<n.
\end{equation}
This implies that \eqref{eq:lem:henk-1} holds with $v=0$ and a nonzero \mbox{$w\in\ker\mathcal{H}_1(x_{[0,T-L]})^\top$}. Next, we consider the case \mbox{$T\geq n+L-1$}. Let $\eta\in\ker \mathcal{H}_{n+L}(u_{[0,T-1]})^\top\backslash\{0\}$.
Define $\eta_0,\ldots,\eta_{n+L-1}\in\mathbb{R}^m$ by $\begin{bmatrix}
\eta_0^\top & \cdots & \eta_{n+L-1}^\top
\end{bmatrix}^\top=\eta$. Observe that for every $t\in[0,T-n-L]$ we have
\begin{equation}
\label{eq:ker2}
\sum_{i=0}^{n+L-1}\eta_i^\top u(t+i)=0.
\end{equation}
Now, we construct a controllable pair $(A,B)$ and an initial condition $x(0)$ such that the generated input-state trajectory satisfies \eqref{eq:lem:henk-1} for some $v\in\mathbb{R}^{mL}$ and $w\in\mathbb{R}^n\backslash\{0\}$. To this end, define
\begin{equation}
\label{eq:Lambda}
\Lambda(\eta)\coloneqq\set{\lambda\in\mathbb{C}}{\textstyle\sum_{i=0}^{n+L-1}\lambda^i \eta_i=0}.
\end{equation}
We note that the set $\Lambda(\eta)$ is finite. Let $A\in\mathbb{R}^{n\times n}$ be a cyclic matrix satisfying\footnote{For example, a Jordan block with its eigenvalue not in $\Lambda(\eta)$.} $\spec A\cap\Lambda(\eta)=\varnothing$. Let $\zeta\in\mathbb{R}^n$ be such that the pair $(A,\zeta)$ is controllable. Define matrices $E_i$, \mbox{$i\in[-1,n+L-1]$}, by the recursion
\begin{equation}
\label{eq:Ei_recursive}
E_{i-1}=AE_{i}+\zeta\eta_i^\top,
\end{equation}
where $E_{n+L-1}=0$. Let $B=E_{-1}=\sum_{i=0}^{n+L-1}A^i\zeta\eta_i^\top$. We claim that $(A,B)$ is controllable. To see this, let\linebreak $z\in\mathbb{C}^n\backslash \{0\}$ be a left eigenvector of $A$, i.e., $z^*A=\lambda z^*$ for some $\lambda\in\spec A$. By the Hautus test and controllability of the pair $(A,\zeta)$, this implies that $z^*\zeta\neq 0$. This in turn implies that $z^* B=z^* \zeta\sum_{i=0}^{n+L-1} \lambda^i \eta_i^\top\neq 0$, thus the pair $(A,B)$ is controllable. Take
\begin{equation}
\label{eq:x(0)}
x(0)=-\sum_{i=0}^{n+L-2} E_i u(i),
\end{equation}
and consider the state trajectory of $(A,B)$, starting from $x(0)$, given by
\begin{equation}
\label{eq:state_construction}
x(t+1)=Ax(t)+Bu(t),\ \hspace{0.25 cm} t\in[0,T-L].
\end{equation}
We claim that this trajectory obeys the following formula with $s\coloneqq T-L-n+1$:
\begin{subequations}
\label{eq:statetraj}
\begin{align}
\label{eq:statetraj-1}
\hspace{-0.75 cm} x(t)&=-\sum_{i=0}^{n+L-2} E_i u(t+i)\hspace{0.25 cm} \text{if}\hspace{0.25 cm} t\in[0,s], \\
\hspace{-0.75 cm} x(t+s)&=\sum_{j=0}^{t-1}\sum_{i=j}^{n+L-2} A^{t-j-1}\zeta\eta_{i-j}^\top u(i+s)\\
\label{eq:statetraj-2}
&-\sum_{i=t}^{n+L-2}E_{i-t} u(i+s)  \hspace{0.25 cm} \text{if}\hspace{0.25 cm} t\in[1,n-1].
\end{align}
\end{subequations}
We use induction to prove \eqref{eq:statetraj-1} and \eqref{eq:statetraj-2}. Note that \eqref{eq:statetraj-1} holds for $t=0$ in view of \eqref{eq:x(0)}. Assume that \eqref{eq:statetraj-1} holds for some $t\in[0,s-1]$. Then, using \eqref{eq:state_construction} it follows that
\begin{subequations}
\mathtoolsset{showonlyrefs=false}
\label{eq:moved}
\begin{align}
&x(t+1)=-A\sum_{i=0}^{n+L-2} E_i u(t+i)+E_{-1}u(t) \\
& \stackrel{\eqref{eq:Ei_recursive}}{=} -\sum_{i=1}^{n+L-2} E_{i-1} u(t+i)+\sum_{i=0}^{n+L-2} \zeta\eta_{i}^\top u(t+i).
\end{align}
\mathtoolsset{showonlyrefs=true}
\end{subequations}
From \eqref{eq:moved} we have
\begin{equation}
\begin{split}
x(t+1)&\stackrel{\eqref{eq:ker2}}{=} -\!\!\!\sum_{i=0}^{n+L-3} \!\! E_{i} u(t\!+\!1\!+\!i)-\zeta\eta_{n+L-1}^\top u(t\!+\!n\!+\!L\!-\!1) \\
&\stackrel{\eqref{eq:Ei_recursive}}{=} -\sum_{i=0}^{n+L-2} E_{i} u(t+1+i),
\end{split}
\end{equation}
which implies that \eqref{eq:statetraj-1} is satisfied for $t+1$. This proves \eqref{eq:statetraj-1}. 

Next, we prove \eqref{eq:statetraj-2}. First, note that \eqref{eq:moved} holds for $t=s$ due to \eqref{eq:state_construction} and \eqref{eq:statetraj-1}. Thus, \eqref{eq:statetraj-2} holds for $t=1$. Assume that \eqref{eq:statetraj-2} holds for some $t\in[1,n-2]$. We observe that using \eqref{eq:state_construction}, we have
\begin{equation}
\label{eq:sth}
\begin{split}
&x(t+s+1)=-\sum_{i=t}^{n+L-2}AE_{i-t} u(s+i) \\ &+\sum_{j=0}^{t-1}\sum_{i=j}^{n+L-2} A^{t-j}\zeta\eta_{i-j}^\top u(s+i)+E_{-1}u(t+s).
\end{split}
\end{equation}
Using \eqref{eq:Ei_recursive}, \eqref{eq:sth} can be reduced to
\begin{equation}
\begin{split}
&x(t+s+1)=-\sum_{i=t+1}^{n+L-2}E_{i-t-1} u(s+i) \\ &+\sum_{i=t}^{n+L-2}\zeta\eta_{i-t} u(s+i)+\sum_{j=0}^{t-1}\sum_{i=j}^{n+L-2} A^{t-j}\zeta\eta_{i-j}^\top u(s+i) \\
&=\!-\!\!\sum_{i=t+1}^{n+L-2}\!E_{i-t-1} u(s\!+\!i)+\sum_{j=0}^{t}\sum_{i=j}^{n+L-2}\! A^{t-j}\zeta\eta_{i-j}^\top u(s\!+\!i),
\end{split}
\end{equation}
which implies that \eqref{eq:statetraj-2} is satisfied for $t+1$. This proves that \eqref{eq:statetraj-2} holds. Now, let $\xi\neq 0$ be such that $\xi^\top A^i \zeta=0$ for all $i\in[0,n-2]$. Observe from \eqref{eq:Ei_recursive} that for every $i\in[L,n+L-1]$ we have $\xi^\top E_i=0$. Thus, we have
\begin{equation}
\label{eq:final_arg1}
\xi^\top\sum_{i=0}^{n+L-2} E_i u(t+i)=\xi^\top\sum_{i=0}^{L-1} E_i u(t+i),
\end{equation}
and for every $t\in[1,n-1]$ we have
\begin{equation}
\label{eq:final_arg2}
\xi^\top\sum_{i=t}^{n+L-2}E_{i-t} u(i+s)=\xi^\top\sum_{i=0}^{L-1} E_i u(i+t+s)
\end{equation}
\begin{equation}
\label{eq:final_arg3}
\text{and }\ \xi^\top\sum_{j=0}^{t-1}\sum_{i=j}^{n+L-2} \!\!A^{t-j-1}\zeta\eta_{i-j}^\top u(i+s)=0.
\end{equation}
Based on \eqref{eq:statetraj-1} and \eqref{eq:statetraj-2}, one can verify using \eqref{eq:final_arg1}, \eqref{eq:final_arg2}, and \eqref{eq:final_arg3} that for every $t\in[0,T-L]$ we have \mbox{$\xi^\top x(t)=-\xi^\top \sum_{i=0}^{L-1} E_i u(t+i)$}. Therefore, \eqref{eq:lem:henk-1} holds with $v=\begin{bmatrix}
E_0 & \cdots & E_{L-1}
\end{bmatrix}^\top \xi$ and $w=\xi$.
\end{proof}

\end{document}
